\documentclass{article}
\usepackage{amsthm}
\usepackage{mathpartir}
\usepackage{url}
\usepackage{imprimatur}
\newtheorem{theorem}{Theorem}
% The paper's environments, as imprimatur.toml's [aliases] declares them.
\leanblockalias{leangrammar}{}
\leanblockalias{leanrules}{}
\leanblockalias{leantheorem}{theorem}
\leanblockalias{leanconventions}{}
\leanblockalias{leannotation}{}
% Conventions of the generated rules: \val is a value, and \rulelabel{Key}{label}
% makes a rule a target for \ref.
\newcommand\val{v}
\makeatletter
\newcommand\rulelabel[2]{\def\@currentlabel{#2}\label{rule:#1}}
\makeatother
\begin{document}

\section{A small calculus, generated from Lean}

This paper presents a simply-typed $\lambda$-calculus whose definitions and
results are in Lean. Everything in a \texttt{lean...} environment below is
generated from the Lean development by imprimatur; the rest is written by
hand, and the hand-written statements are checked against Lean.

\subsection{Syntax}

Terms are name-carrying: a variable is a name, and an abstraction binds a
name. The grammar is generated from the inductive types. The descriptions on
the right are the authors': imprimatur keeps them when it regenerates the
grammar, and uses the constructors' docstrings only for new ones.

\begin{leangrammar}[sum=-gvNYDo1w4]{Stlc.Ty}{Types}
\label{gram:Stlc.Ty}\begin{LeanGrammar}
\LeanProduction{\tau}{types}{\mathsf{b}}{\leanslot{Stlc.Ty.base}{the base type}}
\LeanAlternative{\tau_{1} \to \tau_{2}}{\leanslot{Stlc.Ty.arrow}{functions}}
\end{LeanGrammar}
\end{leangrammar}

\begin{leangrammar}[sum=FuF0yc-Xpz]{Stlc.Tm}{Terms}
\label{gram:Stlc.Tm}\begin{LeanGrammar}
\LeanProduction{e}{terms, with names}{x}{\leanslot{Stlc.Tm.var}{a variable}}
\LeanAlternative{\lambda x{:}\tau.\, e}{\leanslot{Stlc.Tm.lam}{an abstraction, binding $x$ in $e$}}
\LeanAlternative{e_{1}\; e_{2}}{\leanslot{Stlc.Tm.app}{an application}}
\LeanAlternative{()}{\leanslot{Stlc.Tm.unit}{the unit value}}
\end{LeanGrammar}
\end{leangrammar}

\subsection{Names, $\alpha$-equivalence and substitution}

We handle names as Pitts does~\cite{Pitts06}. The basic operation is
\emph{swapping}: $(a\ b) \cdot e$ exchanges the names $a$ and $b$
everywhere in $e$, bound, binding and free occurrences alike. Unlike
renaming, swapping needs no side condition at a binder, and it is its own
inverse; its definition is structural:

\begin{center}
\begin{leanblock}[view=math, def, sum=kQtdL6J5XN]{Stlc.Tm.swap}{Swapping}
\LeanMath{\begin{aligned}(a\ b) \cdot x &= (a\ b)(x) \\ (a\ b) \cdot (\lambda x{:}\tau.\, e) &= \lambda (a\ b)(x){:}\tau.\, (a\ b) \cdot e \\ (a\ b) \cdot (e_{1}\; e_{2}) &= ((a\ b) \cdot e_{1})\; ((a\ b) \cdot e_{2}) \\ (a\ b) \cdot () &= ()\end{aligned}}
\end{leanblock}
\end{center}

Two abstractions are $\alpha$-equivalent when swapping each binder's name
with a name $c$ fresh for both, written $c \mathrel{\#} \ldots$, makes their
bodies $\alpha$-equivalent:

\begin{leanrules}[sum=G1XetiAjYO]{Stlc.Alpha}{Alpha-equivalence}
\begin{mathpar}

\inferrule[($\alpha$-Var)]
  { }
  { x =_{\alpha} x }

\and

\inferrule[($\alpha$-Lam)]
  {
    c \mathrel{\#} a, b, e, e' \\
    (a\ c) \cdot e =_{\alpha} (b\ c) \cdot e'
  }
  { (\lambda a{:}\tau.\, e) =_{\alpha} (\lambda b{:}\tau.\, e') }

\and

\inferrule[($\alpha$-App)]
  {
    e_{1} =_{\alpha} e_{1}' \\
    e_{2} =_{\alpha} e_{2}'
  }
  { e_{1}\; e_{2} =_{\alpha} e_{1}'\; e_{2}' }

\and

\inferrule[($\alpha$-Unit)]
  { }
  { () =_{\alpha} () }

\end{mathpar}
\end{leanrules}

Substitution $e[u/x]$ is then defined by recursion on $e$. At an abstraction
$\lambda a{:}\tau.\, e$ it first swaps $a$ with a name $c$ fresh for $x$,
$a$, $e$ and $u$, so that no free name of $u$ is captured, and continues in
$(a\ c) \cdot e$\leancheck[stmt=a457f283ff65, sig=3e2ddde53b5b]{Stlc.subst}. Swapping preserves the size of a
term\leancheck[stmt=95b8589b3570, sig=c576560f0187]{Stlc.Tm.size_swap}, so the recursion terminates, although
$(a\ c) \cdot e$ is not a subterm of the abstraction.

\subsection{Typing}

The typing rules are generated from the constructors of \texttt{HasType}.
Their labels come from the constructors' docstrings. The well-formedness side
conditions of the Lean rules are left implicit, by a convention declared in
the notation module.

\begin{leanrules}[sum=nDDgEQ1-Ha]{Stlc.HasType}{Typing}
\begin{mathpar}

\inferrule[T-Var]
  {
    \Gamma(x) = \tau
  }
  { \Gamma \vdash x : \tau }

\and

\inferrule[T-Abs]
  {
    \Gamma, x{:}\tau_{1} \vdash e : \tau_{2}
  }
  { \Gamma \vdash (\lambda x{:}\tau_{1}.\, e) : \tau_{1} \to \tau_{2} }

\and

\inferrule[T-App]
  {
    \Gamma \vdash e_{1} : \tau_{1} \to \tau_{2} \\
    \Gamma \vdash e_{2} : \tau_{1}
  }
  { \Gamma \vdash e_{1}\; e_{2} : \tau_{2} }

\and

\inferrule[T-Unit]
  { }
  { \Gamma \vdash () : \mathsf{b} }

\end{mathpar}
\end{leanrules}

\subsection{Reduction}

Reduction is call by value. A premise that a term is a value is left
implicit: the value is written $\val$ instead. Each rule is labelled from its
constructor's name, and can be referred to: rule~\ref{rule:AppLeft}
reduces the function of an application.

\begin{leanrules}[sum=V5YXLse5tR]{Stlc.Step}{Reduction}
\begin{mathpar}

\rulelabel{Beta}{($\longrightarrow$-Beta)}%
\inferrule[($\longrightarrow$-Beta)]
  { }
  { (\lambda x{:}\tau.\, e)\; \val \longrightarrow e[\val/x] }

\and

\rulelabel{AppLeft}{($\longrightarrow$-App Left)}%
\inferrule[($\longrightarrow$-App Left)]
  {
    e_{1} \longrightarrow e_{1}'
  }
  { e_{1}\; e_{2} \longrightarrow e_{1}'\; e_{2} }

\and

\rulelabel{AppRight}{($\longrightarrow$-App Right)}%
\inferrule[($\longrightarrow$-App Right)]
  {
    e_{2} \longrightarrow e_{2}'
  }
  { \val\; e_{2} \longrightarrow \val\; e_{2}' }

\end{mathpar}
\end{leanrules}

\subsection{Results}

A result can be generated, as a sentence; its key can be a title, declared
in the notation module:

\begin{leantheorem}[sum=zzIi1wbp85]{Swapping is involutive}{Swapping is involutive}
\label{thm:Swapping-is-involutive}
\LeanProse{Let $a, b \in \mathit{Name}$ and $e \in \mathit{Term}$. Then $(a\ b) \cdot ((a\ b) \cdot e) = e$.}
\end{leantheorem}

\begin{leantheorem}[sum=OKrVw4uVc8]{Values are final}{Values are final}
\label{thm:Values-are-final}
\LeanProse{Let $\val, e \in \mathit{Term}$. Then $\neg (\val \longrightarrow e)$.}
\end{leantheorem}

Or as a rule, from a theorem whose hypotheses become its premises:
\begin{center}
\begin{leanblock}[view=math, rule, sum=fpyxtM_nmj]{Stlc.step_app}{Congruence}
\LeanMath{\inferrule{e_{1} \longrightarrow e_{1}'}{e_{1}\; e_{2} \longrightarrow e_{1}'\; e_{2}}}
\end{leanblock}
\end{center}

Or the prose can be written by hand and checked: the unit value has only
the base type, that is, if $\Gamma \vdash () : \tau$ then
$\tau = \mathsf{b}$\leancheck[stmt=a39c03ca314b, sig=59daca2b7dd5]{Stlc.unit_type}. Should the Lean statement of
\verb|unit_type| change, \texttt{imprimatur check} reports this sentence.

\subsection{Type safety}

The substitution lemma is proved as Pitts's paper suggests. Typing is
equivariant: swapping two names throughout a term and its context preserves
its type.

\begin{leantheorem}[sum=T-GyIH9gSH]{Equivariance}{Equivariance}
\label{thm:Equivariance}
\LeanProse{Let $a, b \in \mathit{Name}$, $\Gamma \in \mathit{Context}$, $e \in \mathit{Term}$ and $\tau \in \mathit{Type}$. Assume $\Gamma \vdash e : \tau$. Then $(a\ b) \cdot \Gamma \vdash (a\ b) \cdot e : \tau$.}
\end{leantheorem}

Typing also depends on the context only at the term's free names, which
gives weakening and shadowing at once\leancheck[stmt=8f91408eb45c, sig=a65d051214fc]{Stlc.HasType.agree}. So
substitution preserves types: if $\Gamma, x{:}\tau_1 \vdash e : \tau$ and
$\Gamma \vdash u : \tau_1$, then $\Gamma \vdash e[u/x] : \tau$. At an
abstraction, substitution renames the binder to a fresh name by swapping, and
equivariance carries the typing of its body across the renaming; the
induction is on the size of $e$, which swapping
preserves\leancheck[stmt=6842915c5069, sig=105ed38d10fd]{Stlc.subst_typing}. Type safety follows.

\begin{leantheorem}[sum=VQLUk38uFC]{Preservation}{Preservation}
\label{thm:Preservation}
\LeanProse{Let $e, e' \in \mathit{Term}$, $\Gamma \in \mathit{Context}$ and $\tau \in \mathit{Type}$. Assume $e \longrightarrow e'$ and $\Gamma \vdash e : \tau$. Then $\Gamma \vdash e' : \tau$.}
\end{leantheorem}

\begin{leantheorem}[sum=CxgvvaS6H8]{Progress}{Progress}
\label{thm:Progress}
\LeanProse{Let $e \in \mathit{Term}$ and $\tau \in \mathit{Type}$. Assume $\vdash e : \tau$. Then $e \in \mathit{Val} \vee (\exists e',\ e \longrightarrow e')$.}
\end{leantheorem}

\begin{thebibliography}{1}
\bibitem[Pitts(2006)]{Pitts06}
Andrew~M. Pitts.
\newblock Alpha-structural recursion and induction.
\newblock \emph{Journal of the ACM}, 53(3):459--506, 2006.
\newblock \url{https://doi.org/10.1145/1147954.1147961}.
\end{thebibliography}

\appendix
\section{Index of notation}\label{app:notation}

Metavariables range over the types below; a variable with a subscript or a
prime ranges over the same type. imprimatur generates the table from the
categories the notation declares, and checks that every variable printed with
a metavariable has its type.

\begin{leanconventions}[sum=hkPr-RgjND]{Conventions}{Metavariables and the types they range over}
\begin{longtable}{@{}l@{\quad}l@{\quad}p{.4\linewidth}@{\quad}l@{}}
\textit{Variable} & \textit{Type} & & \textit{Shown in}\\
\hline\endhead
$\tau$ & $\mathit{Type}$ & \leanslot{Stlc.Ty}{types} & \textit{Ty}, p.~\pageref{gram:Stlc.Ty}\\
$e$ & $\mathit{Term}$ & \leanslot{Stlc.Tm}{terms, with names} & \textit{Tm}, p.~\pageref{gram:Stlc.Tm}\\
\end{longtable}
A variable with a prime or a subscript ($\tau'$, $\tau_{1}$) ranges over the same type.
\end{leanconventions}

The notation, grouped by where it is defined: the paper's own (\textit{Stlc})
and Lean's. Each is written as the paper writes it, with the Lean constant it
stands for.

\begin{leannotation}[sum=LBK-Px3c45]{Notation}{The notation of the paper}
\begin{longtable}{@{}l@{\quad}l@{\quad}p{.45\linewidth}@{}}
\multicolumn{3}{@{}l}{\textit{Lean}}\\
$\Gamma, x{:}\tau$ & \texttt{List.cons} & \leanslot{List.cons/0}{The list whose first element is \LeanCodeInline{head}, where \LeanCodeInline{tail} is the rest of the list}\\
\multicolumn{3}{@{}l}{\textit{Stlc}}\\
$e =_{\alpha} e'$ & \texttt{Stlc.Alpha} & \leanslot{Stlc.Alpha/0}{α-equivalence}\\
$\mathit{Name}$ & \texttt{Stlc.Atom} & \leanslot{Stlc.Atom/0}{Names of variables}\\
$\mathit{Context}$ & \texttt{Stlc.Ctx} & \leanslot{Stlc.Ctx/0}{Contexts: names with their types, the most recent first}\\
$(a\ b) \cdot \Gamma$ & \texttt{Stlc.Ctx.swap} & \leanslot{Stlc.Ctx.swap/0}{Swapping two names in a context}\\
$\Gamma \vdash e : \tau$ & \texttt{Stlc.HasType} & \leanslot{Stlc.HasType/0}{Typing}\\
$\vdash e : \tau$ & \texttt{Stlc.HasType} & \leanslot{Stlc.HasType/1}{Typing}\\
$v \in \mathit{Val}$ & \texttt{Stlc.IsValue} & \leanslot{Stlc.IsValue/0}{Values}\\
$e \longrightarrow e'$ & \texttt{Stlc.Step} & \leanslot{Stlc.Step/0}{Small-step reduction, call by value}\\
$e_{1}\; e_{2}$ & \texttt{Stlc.Tm.app} & \leanslot{Stlc.Tm.app/0}{an application}\\
$c \mathrel{\#} a, b, e, e'$ & \texttt{Stlc.Tm.atoms} & \leanslot{Stlc.Tm.atoms/0}{The names that occur in a term, bound or free}\\
$\lambda x{:}\tau.\, e$ & \texttt{Stlc.Tm.lam} & \leanslot{Stlc.Tm.lam/0}{an abstraction, binding $x$ in $e$}\\
$(a\ b) \cdot e$ & \texttt{Stlc.Tm.swap} & \leanslot{Stlc.Tm.swap/0}{Swapping two names throughout a term: bound, binding and free occurrences alike}\\
$()$ & \texttt{Stlc.Tm.unit} & \leanslot{Stlc.Tm.unit/0}{the unit value}\\
$x$ & \texttt{Stlc.Tm.var} & \leanslot{Stlc.Tm.var/0}{a variable}\\
$\tau_{1} \to \tau_{2}$ & \texttt{Stlc.Ty.arrow} & \leanslot{Stlc.Ty.arrow/0}{functions}\\
$\mathsf{b}$ & \texttt{Stlc.Ty.base} & \leanslot{Stlc.Ty.base/0}{the base type}\\
$\Gamma(x) = \tau$ & \texttt{Stlc.lookup} & \leanslot{Stlc.lookup/0}{The type of a name in a context}\\
$e[u/x]$ & \texttt{Stlc.subst} & \leanslot{Stlc.subst/0}{Capture-avoiding substitution of \LeanCodeInline{u} for \LeanCodeInline{x} in \LeanCodeInline{e}}\\
$(a\ b)(c)$ & \texttt{Stlc.swapAtom} & \leanslot{Stlc.swapAtom/0}{Swapping two names: \LeanCodeInline{a} becomes \LeanCodeInline{b}, \LeanCodeInline{b} becomes \LeanCodeInline{a}, and other names are unchanged}\\
\end{longtable}
\end{leannotation}

\end{document}
